% Detailed bridge between Born memory kernels and the Gaussian/Langevin form.

\begin{derivation}[从\eqnrefs{eq:born_nonmarkov_redfield_kernel}到\eqnrefs{eq:born_nonmarkov_noise_dissipation_form}]
记

\begin{equation*}
A=q(t),\qquad B_\tau=q(t-\tau),\qquad \rho_\tau=\rho_{S,I}(t-\tau),
\end{equation*}
并把平稳环境相关函数写成
\begin{equation*}
C(\tau)=C_{\rm sym}(\tau)+\ii C_{\rm asym}(\tau),\qquad
C(-\tau)=C_{\rm sym}(\tau)-\ii C_{\rm asym}(\tau).
\end{equation*}
二阶 Born 核中的算符组合为
\begin{align*}
\mathcal K_\tau
={}&C(\tau)[A,B_\tau\rho_\tau]
+C(-\tau)[\rho_\tau B_\tau,A]\\
={}&(C_{\rm sym}+\ii C_{\rm asym})(AB_\tau\rho_\tau-B_\tau\rho_\tau A)
+(C_{\rm sym}-\ii C_{\rm asym})(\rho_\tau B_\tau A-A\rho_\tau B_\tau).
\end{align*}
先收集 $C_{\rm sym}$ 项：
\begin{align*}
\mathcal K_\tau^{({\rm sym})}
&=C_{\rm sym}(AB_\tau\rho_\tau-A\rho_\tau B_\tau-B_\tau\rho_\tau A+\rho_\tau B_\tau A)\\
&=C_{\rm sym}[A,[B_\tau,\rho_\tau]].
\end{align*}
再收集 $C_{\rm asym}$ 项：
\begin{align*}
\mathcal K_\tau^{({\rm asym})}
&=\ii C_{\rm asym}(AB_\tau\rho_\tau+A\rho_\tau B_\tau-B_\tau\rho_\tau A-\rho_\tau B_\tau A)\\
&=\ii C_{\rm asym}[A,\{B_\tau,\rho_\tau\}].
\end{align*}
因此
\begin{equation*}
\mathcal K_\tau
=C_{\rm sym}(\tau)[q(t),[q(t-\tau),\rho_\tau]]
+\ii C_{\rm asym}(\tau)[q(t),\{q(t-\tau),\rho_\tau\}].
\end{equation*}
把它代回二阶 Born 方程前面的 $-\lambda^2/\hbar^2$ 系数和 $\tau$ 积分，就得到\eqnrefs{eq:born_nonmarkov_noise_dissipation_form}。
\end{derivation}

\begin{derivation}[从\eqnrefs{eq:generalized_langevin_from_influence}到\eqnrefs{eq:dressed_oscillator_retarded_response}]
对谐振子势 $V(q)=M\Omega_0^2q^2/2$，正文\eqnrefs{eq:generalized_langevin_from_influence} 可写成
\[
M\ddot q(t)+M\Omega_0^2q(t)
+\int_{-\infty}^{t}\dd s\,\chi_B^R(t-s)q(s)
=\xi(t).
\]
在推迟核中令 $\tau=t-s$，则
\[
\int_{-\infty}^{t}\dd s\,\chi_B^R(t-s)q(s)
=\int_0^\infty\dd\tau\,\chi_B^R(\tau)q(t-\tau).
\]
采用正文的 Fourier 约定
$q(t)=\int \dd\omega\,q(\omega)\ee^{-\ii\omega t}/(2\pi)$，卷积项变成
\begin{align*}
&\int_0^\infty\dd\tau\,\chi_B^R(\tau)
\int\frac{\dd\omega}{2\pi}q(\omega)
\ee^{-\ii\omega(t-\tau)}\\
&=\int\frac{\dd\omega}{2\pi}
\left[
\int_0^\infty\dd\tau\,
\chi_B^R(\tau)\ee^{+\ii\omega\tau}
\right]q(\omega)\ee^{-\ii\omega t}\\
&=\int\frac{\dd\omega}{2\pi}
\chi_B^R(\omega)q(\omega)\ee^{-\ii\omega t}.
\end{align*}
同时 $\ddot q(t)\mapsto-\omega^2q(\omega)$。因此每个频率分量满足
\[
\left[
M\bigl(\Omega_0^2-(\omega+\ii0^+)^2\bigr)
+\chi_B^R(\omega)
\right]q(\omega)=\xi(\omega).
\]
对乘在 $q(\omega)$ 前的标量核求逆，得到
\[
q(\omega)=\mathcal R_q(\omega)\xi(\omega),
\qquad
\mathcal R_q(\omega)
=\frac{1}{M[\Omega_0^2-(\omega+\ii0^+)^2]+\chi_B^R(\omega)},
\]
即正文\eqnrefs{eq:dressed_oscillator_retarded_response}。
\end{derivation}

\begin{derivation}[从\eqnrefs{eq:dressed_coordinate_spectrum_dp68}到\eqnrefs{eq:dressed_system_fdt_from_bath_fdt}]

正文\eqnrefs{eq:dressed_coordinate_spectrum_dp68} 表明系统坐标谱等于响应模方乘环境噪声谱。热平衡时，环境噪声谱又由\eqnrefs{eq:fdt_from_defined_correlators} 与环境响应的虚部关联。利用复数逆函数的虚部恒等式，把 $|\mathcal R_{x_S}|^2[-\operatorname{Im}\chi_B^R]$ 改写成系统自身响应的耗散部分，就得到\eqnrefs{eq:dressed_system_fdt_from_bath_fdt}。

令
\[
D(\omega)=M[\Omega_0^2-(\omega+\ii0^+)^2]+\chi_B^R(\omega),
\qquad
\mathcal R_{x_S}=D^{-1}.
\]
于是
\[
\operatorname{Im}\mathcal R_{x_S}
=-|\mathcal R_{x_S}|^2\operatorname{Im}\chi_B^R.
\]
按正文约定 $\chi_{x_S}^R=-\mathcal R_{x_S}$，因此
\[
-\operatorname{Im}\chi_{x_S}^R
=|\mathcal R_{x_S}|^2[-\operatorname{Im}\chi_B^R].
\]
将这一恒等式和环境 FDT~\eqnrefs{eq:fdt_from_defined_correlators} 代入\eqnrefs{eq:dressed_coordinate_spectrum_dp68}，得到\eqnrefs{eq:dressed_system_fdt_from_bath_fdt}。
\end{derivation}
